
\subsection{The Bound}

\subsubsection{Projectors}
Start with the basics, we follow Trefethen and Bau. A projector
/idempotent is a square matrix $P^2 = P$. A projector separates
$\mathbb{R}^m$ into two spaces. An orthogonal projctor is a porjector
that is symmetric $P^t = P$. We can also write $P= QQ^t$ where the
columns of $Q$ are orthogonal. If a subspace is spanned by vectors
$\{a_1,\dots,a_n\}$ we take $A$ to be the matrix whose columns are
$a_j$. The \textit{orthogonal} projection $y$ of an arbitrary vector
$v$ onto $\text{range}(A)$ must satistfy $y-v \perp \text{range}(A)$,
equivalently $a_j^t(y-v) = 0$. Since $y \in \text{range}(A)$ we can
take $y =Ax$ and the above condition becomes $a_j^t(Ax-v) = 0$, or
$A^t(Ax-v) = 0$, or $A^tAx = A^tv$. Since $A$ is full rank, $A^tA$ is
invertible and then $x = (A^tA)^{-1}A^tv$. This leads to the known
formula of the orthogonal projector of an arbitrary $v$ onto the
column space of $A$ as $y = A(A^tA)^{-1}A^tv$. The orthogonal
projector onto $\text{range}(A)$ is thus

$$
P = A(A^tA)^{-1}A^t
$$


\subsubsection{Conditioning}
We defin the matrix norm as $\|A\| = \sup_{x\neq 0}
\frac{\|Ax\|}{\|x\|} = \sup_{\|x\|=1} \|Ax\|$. For a ssquare matrix,
this is the maximal absolute eigenvalue of $A$. We define the
condition number of a matrix as $\kappa(A) = \|A\| \|A^{-1}\|$. For a
square matrix, this is the ratio of maximal to minimal absolute
eigenvalues.



\section{Vector Autoregression (VAR) Analysis}

\subsection{Model Framework}

To empirically test the feasibility of connectivity inference, we employ Vector Autoregression models on real pneumonia and influenza incidence data from US states.

Let $\mathbf{X}_t = (X_t^{(1)}, X_t^{(2)}, \ldots, X_t^{(n)})^T$ denote the vector of incidence rates across $n$ states at time $t$, where $t = 1, 2, \ldots, T$ with $T = 52 \times 30$ (30 years of weekly data).

\subsection{VAR Model Specifications}

\subsubsection{Model 1: Full VAR (Connectivity Allowed)}
The unrestricted VAR(p) model allows for full connectivity between states:
\begin{equation}
\mathbf{X}_t = \mathbf{A}_1 \mathbf{X}_{t-1} + \mathbf{A}_2 \mathbf{X}_{t-2} + \cdots + \mathbf{A}_p \mathbf{X}_{t-p} + {\varepsilon}_t \label{eq:var_full}
\end{equation}

where $\mathbf{A}_i \in \mathbb{R}^{n \times n}$ are coefficient matrices with unrestricted entries, and ${\varepsilon}_t \sim \mathcal{N}(\mathbf{0}, {\Sigma})$ is white noise.

The off-diagonal elements of $\mathbf{A}_i$ capture cross-state dependencies, representing potential connectivity effects.

\subsubsection{Model 2: Diagonal VAR (No Connectivity)}
The restricted model constrains coefficient matrices to be diagonal:
\begin{equation}
\mathbf{X}_t = \mathbf{D}_1 \mathbf{X}_{t-1} + \mathbf{D}_2 \mathbf{X}_{t-2} + \cdots + \mathbf{D}_p \mathbf{X}_{t-p} + {\varepsilon}_t \label{eq:var_diagonal}
\end{equation}

where $\mathbf{D}_i = \text{diag}(d_{i1}, d_{i2}, \ldots, d_{in})$ are diagonal matrices. This model is equivalent to fitting independent ARIMA models to each state.

\subsubsection{Model 3: Univariate SARIMA Ensemble}
For comparison, we fit individual Seasonal ARIMA models to each state:
\begin{equation}
\phi_i(L)(1-L)^{d_i}(1-L^s)^{D_i} X_t^{(i)} = \theta_i(L)(1-L^s)^{D_i}\varepsilon_t^{(i)} \label{eq:sarima}
\end{equation}

where $L$ is the lag operator, $s=52$ is the seasonal period, and $\phi_i(L)$, $\theta_i(L)$ are lag polynomials.

\subsection{Model Selection and Estimation}

\subsubsection{Lag Order Selection}
Optimal lag order $p$ is selected using information criteria:
\begin{align}
\text{AIC}(p) &= \log|\hat{{\Sigma}}(p)| + \frac{2k}{T} \label{eq:aic}\\
\text{BIC}(p) &= \log|\hat{{\Sigma}}(p)| + \frac{k \log T}{T} \label{eq:bic}
\end{align}

where $k = n^2 p$ is the number of parameters and $\hat{{\Sigma}}(p)$ is the estimated error covariance matrix.

\subsubsection{Estimation Method}
VAR models are estimated using ordinary least squares (OLS) applied equation-by-equation, which is efficient and consistent under standard conditions.

\subsection{Model Comparison Framework}

\subsubsection{Information Criteria}
Models are compared using:
\begin{itemize}
\item Akaike Information Criterion (AIC)
\item Bayesian Information Criterion (BIC)  
\item Hannan-Quinn Information Criterion (HQIC)
\end{itemize}

\subsubsection{Forecasting Performance}
Out-of-sample forecasting accuracy is evaluated using:
\begin{align}
\text{RMSE} &= \sqrt{\frac{1}{h}\sum_{j=1}^{h} \|\mathbf{X}_{T+j} - \hat{\mathbf{X}}_{T+j|T}\|^2} \label{eq:rmse}\\
\text{MAE} &= \frac{1}{h}\sum_{j=1}^{h} \|\mathbf{X}_{T+j} - \hat{\mathbf{X}}_{T+j|T}\|_1 \label{eq:mae}
\end{align}

where $h$ is the forecast horizon.

\subsubsection{Connectivity Significance Tests}
To test whether connectivity terms significantly improve model performance, we employ:

1. \textbf{Likelihood Ratio Tests}: Compare nested models (diagonal vs. full VAR)
2. \textbf{Granger Causality Tests}: Test whether lagged values of state $j$ help predict state $i$
3. \textbf{Cross-validation}: Compare out-of-sample forecasting performance

\subsection{Stationarity and Preprocessing}

\subsubsection{Unit Root Tests}
Data stationarity is assessed using:
\begin{itemize}
\item Augmented Dickey-Fuller (ADF) tests for each series
\item KPSS tests for stationarity confirmation
\item Johansen cointegration tests for long-run relationships
\end{itemize}

\subsubsection{Data Transformation}
Non-stationary series are transformed via:
\begin{equation}
\tilde{X}_t^{(i)} = \Delta^{d_i} \log(X_t^{(i)} + c) \label{eq:transform}
\end{equation}

where $\Delta^{d_i}$ denotes $d_i$-th order differencing and $c$ is a small constant to handle zero values.

\section{Integration of Theoretical and Empirical Results}

The SEIR model analysis provides theoretical predictions about the ill-conditioning of connectivity inference, quantified by condition numbers $\kappa(\mathbf{M})$. The VAR analysis provides empirical validation using real surveillance data.

The key hypothesis is that if connectivity inference is fundamentally impossible due to epidemic synchronization (as predicted by theory), then:

1. Full VAR models should not significantly outperform diagonal VAR models
2. Off-diagonal coefficient estimates should be unstable and statistically insignificant  
3. Out-of-sample forecasting performance should favor simpler (no-connectivity) models

This framework provides a comprehensive test of the practical feasibility of connectivity inference from epidemic surveillance data.

% End of methods section
